\hypertarget{p1-building-the-ultimate-tic-tac-toe-engine-in-python}{%
\section{P1: Building the Ultimate Tic-Tac-Toe Engine in
Python}\label{p1-building-the-ultimate-tic-tac-toe-engine-in-python}}

\emph{This is the first practitioner post in the series. The theory
essays explain why things work; the practitioner posts show how to build
them. You need Python 3.9+ and NumPy. All code in this post is
self-contained --- you can paste it directly into a notebook or script.
The full runnable notebook is available on
\href{https://colab.research.google.com/github/thltsui/UlltimateTicTacToe/blob/Substack/substack/notebook_1_uttt_engine.ipynb}{Google
Colab}.}

\begin{center}\rule{0.5\linewidth}{0.5pt}\end{center}

Before we can train any AI, we need a game engine: something that knows
the rules, generates legal moves, applies them, and tells us when the
game is over. This post builds that engine from scratch. By the end, you
will have a fully functional UTTT implementation and an understanding of
how game state is encoded as a tensor for neural network input.

\hypertarget{representing-the-board}{%
\subsection{1. Representing the Board}\label{representing-the-board}}

The first design decision is how to represent the game state. We want
something that is easy to reason about, efficient to copy (we will be
making millions of copies during search), and expressive enough to
capture everything the rules require.

A \texttt{dataclass} with NumPy arrays works well:

\begin{Shaded}
\begin{Highlighting}[]
\ImportTok{import}\NormalTok{ numpy }\ImportTok{as}\NormalTok{ np}
\ImportTok{from}\NormalTok{ dataclasses }\ImportTok{import}\NormalTok{ dataclass, field}

\AttributeTok{@dataclass}
\KeywordTok{class}\NormalTok{ GameState:}
\NormalTok{    cells: np.ndarray }\OperatorTok{=}\NormalTok{ field(}
\NormalTok{        default\_factory}\OperatorTok{=}\KeywordTok{lambda}\NormalTok{: np.zeros((}\DecValTok{9}\NormalTok{, }\DecValTok{9}\NormalTok{), dtype}\OperatorTok{=}\NormalTok{np.int8)}
\NormalTok{    )}
\NormalTok{    sub\_board\_results: np.ndarray }\OperatorTok{=}\NormalTok{ field(}
\NormalTok{        default\_factory}\OperatorTok{=}\KeywordTok{lambda}\NormalTok{: np.zeros(}\DecValTok{9}\NormalTok{, dtype}\OperatorTok{=}\NormalTok{np.int8)}
\NormalTok{    )}
\NormalTok{    active\_sub\_board: }\BuiltInTok{int} \OperatorTok{=} \OperatorTok{{-}}\DecValTok{1}
\NormalTok{    current\_player: }\BuiltInTok{int} \OperatorTok{=} \DecValTok{1}
\NormalTok{    move\_count: }\BuiltInTok{int} \OperatorTok{=} \DecValTok{0}
\NormalTok{    is\_terminal: }\BuiltInTok{bool} \OperatorTok{=} \VariableTok{False}
\NormalTok{    winner: }\BuiltInTok{object} \OperatorTok{=} \VariableTok{None}
\end{Highlighting}
\end{Shaded}

Breaking this down field by field:

\textbf{\texttt{cells} --- shape (9, 9):} The 81 squares of the board,
indexed as \texttt{cells{[}sub\_board,\ cell{]}}. Values are \texttt{+1}
(X), \texttt{-1} (O), or \texttt{0} (empty). Sub-boards and cells are
both indexed 0--8, left-to-right then top-to-bottom within each 3×3
grid.

\textbf{\texttt{sub\_board\_results} --- shape (9,):} The status of each
of the 9 local boards. Values are \texttt{+1} (X won), \texttt{-1} (O
won), \texttt{2} (draw), or \texttt{0} (still in play).

\textbf{\texttt{active\_sub\_board}:} Which local board the current
player must play in. \texttt{-1} means free choice --- the player may
play anywhere.

\textbf{\texttt{current\_player}:} \texttt{+1} for X, \texttt{-1} for O.
This sign convention makes win detection and symmetry handling elegant.

\textbf{\texttt{is\_terminal} and \texttt{winner}:} Terminal states end
the game. \texttt{winner} is \texttt{+1} (X wins), \texttt{-1} (O wins),
or \texttt{0} (draw).

\hypertarget{why-immutability-matters}{%
\subsubsection{Why immutability
matters}\label{why-immutability-matters}}

We make the game state immutable by always creating a copy before
applying a move:

\begin{Shaded}
\begin{Highlighting}[]
\KeywordTok{def}\NormalTok{ copy(}\VariableTok{self}\NormalTok{):}
    \ControlFlowTok{return}\NormalTok{ GameState(}
\NormalTok{        cells}\OperatorTok{=}\VariableTok{self}\NormalTok{.cells.copy(),}
\NormalTok{        sub\_board\_results}\OperatorTok{=}\VariableTok{self}\NormalTok{.sub\_board\_results.copy(),}
\NormalTok{        active\_sub\_board}\OperatorTok{=}\VariableTok{self}\NormalTok{.active\_sub\_board,}
\NormalTok{        current\_player}\OperatorTok{=}\VariableTok{self}\NormalTok{.current\_player,}
\NormalTok{        move\_count}\OperatorTok{=}\VariableTok{self}\NormalTok{.move\_count,}
\NormalTok{        is\_terminal}\OperatorTok{=}\VariableTok{self}\NormalTok{.is\_terminal,}
\NormalTok{        winner}\OperatorTok{=}\VariableTok{self}\NormalTok{.winner,}
\NormalTok{    )}
\end{Highlighting}
\end{Shaded}

Immutable states mean that the tree search can hold references to many
board positions simultaneously without any of them corrupting each
other. Mutable state is a source of hard-to-find bugs in game engines.

\hypertarget{move-encoding}{%
\subsection{2. Move Encoding}\label{move-encoding}}

UTTT has 81 possible squares. Rather than representing a move as a
\texttt{(sub\_board,\ cell)} pair, we encode each move as a single
integer from 0 to 80:

\begin{Shaded}
\begin{Highlighting}[]
\KeywordTok{def}\NormalTok{ encode\_move(sub\_board: }\BuiltInTok{int}\NormalTok{, cell: }\BuiltInTok{int}\NormalTok{) }\OperatorTok{{-}\textgreater{}} \BuiltInTok{int}\NormalTok{:}
    \ControlFlowTok{return}\NormalTok{ sub\_board }\OperatorTok{*} \DecValTok{9} \OperatorTok{+}\NormalTok{ cell}

\KeywordTok{def}\NormalTok{ decode\_move(move\_idx: }\BuiltInTok{int}\NormalTok{) }\OperatorTok{{-}\textgreater{}} \BuiltInTok{tuple}\NormalTok{[}\BuiltInTok{int}\NormalTok{, }\BuiltInTok{int}\NormalTok{]:}
    \ControlFlowTok{return}\NormalTok{ move\_idx }\OperatorTok{//} \DecValTok{9}\NormalTok{, move\_idx }\OperatorTok{\%} \DecValTok{9}
\end{Highlighting}
\end{Shaded}

This flat encoding is convenient for the neural network's policy head,
which outputs a probability distribution over all 81 moves. Sub-board 0,
cell 0 is move 0; sub-board 8, cell 8 is move 80.

The spatial position of a cell in the 9×9 grid (for visualisation) is:

\begin{Shaded}
\begin{Highlighting}[]
\NormalTok{sub\_board, cell }\OperatorTok{=}\NormalTok{ decode\_move(move\_idx)}
\NormalTok{row }\OperatorTok{=}\NormalTok{ (sub\_board }\OperatorTok{//} \DecValTok{3}\NormalTok{) }\OperatorTok{*} \DecValTok{3} \OperatorTok{+}\NormalTok{ (cell }\OperatorTok{//} \DecValTok{3}\NormalTok{)}
\NormalTok{col }\OperatorTok{=}\NormalTok{ (sub\_board  }\OperatorTok{\%} \DecValTok{3}\NormalTok{) }\OperatorTok{*} \DecValTok{3} \OperatorTok{+}\NormalTok{ (cell  }\OperatorTok{\%} \DecValTok{3}\NormalTok{)}
\end{Highlighting}
\end{Shaded}

\hypertarget{generating-legal-moves}{%
\subsection{3. Generating Legal Moves}\label{generating-legal-moves}}

Legal move generation is the core of the engine. It implements the
``send your opponent'' rule:

\begin{Shaded}
\begin{Highlighting}[]
\KeywordTok{def}\NormalTok{ get\_legal\_moves(state: GameState) }\OperatorTok{{-}\textgreater{}} \BuiltInTok{list}\NormalTok{[}\BuiltInTok{int}\NormalTok{]:}
    \ControlFlowTok{if}\NormalTok{ state.is\_terminal:}
        \ControlFlowTok{return}\NormalTok{ []}
\NormalTok{    moves }\OperatorTok{=}\NormalTok{ []}
    \ControlFlowTok{if}\NormalTok{ state.active\_sub\_board }\OperatorTok{==} \OperatorTok{{-}}\DecValTok{1}\NormalTok{:}
        \CommentTok{\# Free choice: any empty cell in any undecided sub{-}board}
        \ControlFlowTok{for}\NormalTok{ sb }\KeywordTok{in} \BuiltInTok{range}\NormalTok{(}\DecValTok{9}\NormalTok{):}
            \ControlFlowTok{if}\NormalTok{ state.sub\_board\_results[sb] }\OperatorTok{!=} \DecValTok{0}\NormalTok{:}
                \ControlFlowTok{continue}
            \ControlFlowTok{for}\NormalTok{ cell }\KeywordTok{in} \BuiltInTok{range}\NormalTok{(}\DecValTok{9}\NormalTok{):}
                \ControlFlowTok{if}\NormalTok{ state.cells[sb, cell] }\OperatorTok{==} \DecValTok{0}\NormalTok{:}
\NormalTok{                    moves.append(encode\_move(sb, cell))}
    \ControlFlowTok{else}\NormalTok{:}
\NormalTok{        sb }\OperatorTok{=}\NormalTok{ state.active\_sub\_board}
        \ControlFlowTok{if}\NormalTok{ state.sub\_board\_results[sb] }\OperatorTok{==} \DecValTok{0}\NormalTok{:}
            \CommentTok{\# Normal case: play in the required sub{-}board}
            \ControlFlowTok{for}\NormalTok{ cell }\KeywordTok{in} \BuiltInTok{range}\NormalTok{(}\DecValTok{9}\NormalTok{):}
                \ControlFlowTok{if}\NormalTok{ state.cells[sb, cell] }\OperatorTok{==} \DecValTok{0}\NormalTok{:}
\NormalTok{                    moves.append(encode\_move(sb, cell))}
        \ControlFlowTok{if} \KeywordTok{not}\NormalTok{ moves:}
            \CommentTok{\# The required sub{-}board is full or already decided:}
            \CommentTok{\# fall back to free choice}
            \ControlFlowTok{for}\NormalTok{ sb2 }\KeywordTok{in} \BuiltInTok{range}\NormalTok{(}\DecValTok{9}\NormalTok{):}
                \ControlFlowTok{if}\NormalTok{ state.sub\_board\_results[sb2] }\OperatorTok{!=} \DecValTok{0}\NormalTok{:}
                    \ControlFlowTok{continue}
                \ControlFlowTok{for}\NormalTok{ cell }\KeywordTok{in} \BuiltInTok{range}\NormalTok{(}\DecValTok{9}\NormalTok{):}
                    \ControlFlowTok{if}\NormalTok{ state.cells[sb2, cell] }\OperatorTok{==} \DecValTok{0}\NormalTok{:}
\NormalTok{                        moves.append(encode\_move(sb2, cell))}
    \ControlFlowTok{return}\NormalTok{ moves}
\end{Highlighting}
\end{Shaded}

The key logic: if \texttt{active\_sub\_board} is \texttt{-1}, any empty
cell on any undecided board is legal. Otherwise, the current player must
play in \texttt{active\_sub\_board} --- unless that board is already
decided, in which case they get free choice.

For the neural network, we also need a binary mask over all 81 moves:

\begin{Shaded}
\begin{Highlighting}[]
\KeywordTok{def}\NormalTok{ get\_legal\_move\_mask(state: GameState) }\OperatorTok{{-}\textgreater{}}\NormalTok{ np.ndarray:}
\NormalTok{    mask }\OperatorTok{=}\NormalTok{ np.zeros(}\DecValTok{81}\NormalTok{, dtype}\OperatorTok{=}\NormalTok{np.float32)}
    \ControlFlowTok{for}\NormalTok{ move }\KeywordTok{in}\NormalTok{ get\_legal\_moves(state):}
\NormalTok{        mask[move] }\OperatorTok{=} \FloatTok{1.0}
    \ControlFlowTok{return}\NormalTok{ mask}
\end{Highlighting}
\end{Shaded}

This mask is multiplied with the policy head output during search to
zero out illegal moves before the softmax.

\hypertarget{win-detection}{%
\subsection{4. Win Detection}\label{win-detection}}

Win detection runs at two levels: local (did someone win this
sub-board?) and global (did someone win three sub-boards in a row?).

\begin{Shaded}
\begin{Highlighting}[]
\NormalTok{WIN\_LINES }\OperatorTok{=}\NormalTok{ [}
\NormalTok{    (}\DecValTok{0}\NormalTok{,}\DecValTok{1}\NormalTok{,}\DecValTok{2}\NormalTok{), (}\DecValTok{3}\NormalTok{,}\DecValTok{4}\NormalTok{,}\DecValTok{5}\NormalTok{), (}\DecValTok{6}\NormalTok{,}\DecValTok{7}\NormalTok{,}\DecValTok{8}\NormalTok{),   }\CommentTok{\# rows}
\NormalTok{    (}\DecValTok{0}\NormalTok{,}\DecValTok{3}\NormalTok{,}\DecValTok{6}\NormalTok{), (}\DecValTok{1}\NormalTok{,}\DecValTok{4}\NormalTok{,}\DecValTok{7}\NormalTok{), (}\DecValTok{2}\NormalTok{,}\DecValTok{5}\NormalTok{,}\DecValTok{8}\NormalTok{),   }\CommentTok{\# columns}
\NormalTok{    (}\DecValTok{0}\NormalTok{,}\DecValTok{4}\NormalTok{,}\DecValTok{8}\NormalTok{), (}\DecValTok{2}\NormalTok{,}\DecValTok{4}\NormalTok{,}\DecValTok{6}\NormalTok{),             }\CommentTok{\# diagonals}
\NormalTok{]}

\KeywordTok{def}\NormalTok{ check\_sub\_board\_winner(cells\_for\_sb: np.ndarray) }\OperatorTok{{-}\textgreater{}} \BuiltInTok{int}\NormalTok{:}
    \CommentTok{"""Returns +1, {-}1, 2 (draw), or 0 (ongoing)."""}
    \ControlFlowTok{for}\NormalTok{ a, b, c }\KeywordTok{in}\NormalTok{ WIN\_LINES:}
        \ControlFlowTok{if}\NormalTok{ cells\_for\_sb[a] }\OperatorTok{!=} \DecValTok{0} \KeywordTok{and}\NormalTok{ cells\_for\_sb[a] }\OperatorTok{==}\NormalTok{ cells\_for\_sb[b] }\OperatorTok{==}\NormalTok{ cells\_for\_sb[c]:}
            \ControlFlowTok{return} \BuiltInTok{int}\NormalTok{(cells\_for\_sb[a])}
    \ControlFlowTok{if}\NormalTok{ np.}\BuiltInTok{all}\NormalTok{(cells\_for\_sb }\OperatorTok{!=} \DecValTok{0}\NormalTok{):}
        \ControlFlowTok{return} \DecValTok{2}  \CommentTok{\# draw}
    \ControlFlowTok{return} \DecValTok{0}      \CommentTok{\# still in play}

\KeywordTok{def}\NormalTok{ check\_meta\_winner(sub\_board\_results: np.ndarray) }\OperatorTok{{-}\textgreater{}} \BuiltInTok{int}\NormalTok{:}
    \CommentTok{"""Returns +1, {-}1, 2 (draw), or 0 (ongoing) for the global board."""}
    \ControlFlowTok{for}\NormalTok{ a, b, c }\KeywordTok{in}\NormalTok{ WIN\_LINES:}
        \ControlFlowTok{if}\NormalTok{ (sub\_board\_results[a] }\KeywordTok{in}\NormalTok{ (}\DecValTok{1}\NormalTok{, }\OperatorTok{{-}}\DecValTok{1}\NormalTok{) }\KeywordTok{and}
\NormalTok{                sub\_board\_results[a] }\OperatorTok{==}\NormalTok{ sub\_board\_results[b] }\OperatorTok{==}\NormalTok{ sub\_board\_results[c]):}
            \ControlFlowTok{return} \BuiltInTok{int}\NormalTok{(sub\_board\_results[a])}
    \ControlFlowTok{if}\NormalTok{ np.}\BuiltInTok{all}\NormalTok{(sub\_board\_results }\OperatorTok{!=} \DecValTok{0}\NormalTok{):}
        \ControlFlowTok{return} \DecValTok{2}
    \ControlFlowTok{return} \DecValTok{0}
\end{Highlighting}
\end{Shaded}

\hypertarget{applying-a-move}{%
\subsection{5. Applying a Move}\label{applying-a-move}}

\texttt{apply\_move} is the heart of the engine. It takes a state and a
move index, and returns a new state:

\begin{Shaded}
\begin{Highlighting}[]
\KeywordTok{def}\NormalTok{ apply\_move(state: GameState, move\_idx: }\BuiltInTok{int}\NormalTok{) }\OperatorTok{{-}\textgreater{}}\NormalTok{ GameState:}
\NormalTok{    sb, cell }\OperatorTok{=}\NormalTok{ decode\_move(move\_idx)}
\NormalTok{    s }\OperatorTok{=}\NormalTok{ state.copy()}

    \CommentTok{\# Place the mark}
\NormalTok{    s.cells[sb, cell] }\OperatorTok{=}\NormalTok{ state.current\_player}

    \CommentTok{\# Update sub{-}board result if it just ended}
    \ControlFlowTok{if}\NormalTok{ s.sub\_board\_results[sb] }\OperatorTok{==} \DecValTok{0}\NormalTok{:}
\NormalTok{        result }\OperatorTok{=}\NormalTok{ check\_sub\_board\_winner(s.cells[sb])}
        \ControlFlowTok{if}\NormalTok{ result }\OperatorTok{!=} \DecValTok{0}\NormalTok{:}
\NormalTok{            s.sub\_board\_results[sb] }\OperatorTok{=}\NormalTok{ result}

    \CommentTok{\# Check if the global game is over}
\NormalTok{    meta }\OperatorTok{=}\NormalTok{ check\_meta\_winner(s.sub\_board\_results)}
    \ControlFlowTok{if}\NormalTok{ meta }\OperatorTok{!=} \DecValTok{0}\NormalTok{:}
\NormalTok{        s.is\_terminal }\OperatorTok{=} \VariableTok{True}
\NormalTok{        s.winner }\OperatorTok{=} \DecValTok{0} \ControlFlowTok{if}\NormalTok{ meta }\OperatorTok{==} \DecValTok{2} \ControlFlowTok{else}\NormalTok{ meta}

    \CommentTok{\# The cell played determines where the opponent must go next}
    \CommentTok{\# If that sub{-}board is already decided, give free choice}
\NormalTok{    s.active\_sub\_board }\OperatorTok{=} \OperatorTok{{-}}\DecValTok{1} \ControlFlowTok{if}\NormalTok{ s.sub\_board\_results[cell] }\OperatorTok{!=} \DecValTok{0} \ControlFlowTok{else}\NormalTok{ cell}

    \CommentTok{\# Switch player and increment move count}
\NormalTok{    s.current\_player }\OperatorTok{=} \OperatorTok{{-}}\NormalTok{state.current\_player}
\NormalTok{    s.move\_count }\OperatorTok{=}\NormalTok{ state.move\_count }\OperatorTok{+} \DecValTok{1}

    \ControlFlowTok{return}\NormalTok{ s}
\end{Highlighting}
\end{Shaded}

The line
\texttt{s.active\_sub\_board\ =\ -1\ if\ s.sub\_board\_results{[}cell{]}\ !=\ 0\ else\ cell}
is the entire ``send your opponent'' rule in one expression. The cell
index within the current sub-board (\texttt{cell}) determines the next
active sub-board. If \texttt{sub\_board\_results{[}cell{]}} is non-zero
(that sub-board is decided), we set \texttt{active\_sub\_board\ =\ -1}
to grant free choice.

\hypertarget{encoding-state-as-a-tensor}{%
\subsection{6. Encoding State as a
Tensor}\label{encoding-state-as-a-tensor}}

The game state needs to be converted to a tensor for the neural network.
We use 7 channels, each a 9×9 binary grid, always expressed from the
perspective of the current player (so the network always sees ``my
pieces'' in channel 0, regardless of whether the current player is X or
O):

\begin{Shaded}
\begin{Highlighting}[]
\KeywordTok{def}\NormalTok{ encode\_state(state: GameState) }\OperatorTok{{-}\textgreater{}}\NormalTok{ np.ndarray:}
    \CommentTok{"""Returns shape (7, 9, 9) float32 tensor."""}
\NormalTok{    p }\OperatorTok{=}\NormalTok{ state.current\_player  }\CommentTok{\# +1 or {-}1}
\NormalTok{    tensor }\OperatorTok{=}\NormalTok{ np.zeros((}\DecValTok{7}\NormalTok{, }\DecValTok{9}\NormalTok{, }\DecValTok{9}\NormalTok{), dtype}\OperatorTok{=}\NormalTok{np.float32)}

    \ControlFlowTok{for}\NormalTok{ sb }\KeywordTok{in} \BuiltInTok{range}\NormalTok{(}\DecValTok{9}\NormalTok{):}
\NormalTok{        rb, cb }\OperatorTok{=}\NormalTok{ sb }\OperatorTok{//} \DecValTok{3}\NormalTok{, sb }\OperatorTok{\%} \DecValTok{3}
\NormalTok{        result }\OperatorTok{=}\NormalTok{ state.sub\_board\_results[sb]}

        \ControlFlowTok{for}\NormalTok{ cell }\KeywordTok{in} \BuiltInTok{range}\NormalTok{(}\DecValTok{9}\NormalTok{):}
\NormalTok{            r }\OperatorTok{=}\NormalTok{ rb }\OperatorTok{*} \DecValTok{3} \OperatorTok{+}\NormalTok{ cell }\OperatorTok{//} \DecValTok{3}
\NormalTok{            c }\OperatorTok{=}\NormalTok{ cb }\OperatorTok{*} \DecValTok{3} \OperatorTok{+}\NormalTok{ cell }\OperatorTok{\%} \DecValTok{3}
\NormalTok{            v }\OperatorTok{=}\NormalTok{ state.cells[sb, cell]}

            \CommentTok{\# Channel 0: current player\textquotesingle{}s pieces}
            \ControlFlowTok{if}\NormalTok{ v }\OperatorTok{==}\NormalTok{ p:}
\NormalTok{                tensor[}\DecValTok{0}\NormalTok{, r, c] }\OperatorTok{=} \FloatTok{1.0}
            \CommentTok{\# Channel 1: opponent\textquotesingle{}s pieces}
            \ControlFlowTok{elif}\NormalTok{ v }\OperatorTok{==} \OperatorTok{{-}}\NormalTok{p:}
\NormalTok{                tensor[}\DecValTok{1}\NormalTok{, r, c] }\OperatorTok{=} \FloatTok{1.0}

        \CommentTok{\# Channels 2{-}4: sub{-}board outcomes (current player won, opponent won, draw)}
        \ControlFlowTok{if}\NormalTok{ result }\OperatorTok{==}\NormalTok{ p:}
\NormalTok{            tensor[}\DecValTok{2}\NormalTok{, rb}\OperatorTok{*}\DecValTok{3}\NormalTok{:rb}\OperatorTok{*}\DecValTok{3}\OperatorTok{+}\DecValTok{3}\NormalTok{, cb}\OperatorTok{*}\DecValTok{3}\NormalTok{:cb}\OperatorTok{*}\DecValTok{3}\OperatorTok{+}\DecValTok{3}\NormalTok{] }\OperatorTok{=} \FloatTok{1.0}
        \ControlFlowTok{elif}\NormalTok{ result }\OperatorTok{==} \OperatorTok{{-}}\NormalTok{p:}
\NormalTok{            tensor[}\DecValTok{3}\NormalTok{, rb}\OperatorTok{*}\DecValTok{3}\NormalTok{:rb}\OperatorTok{*}\DecValTok{3}\OperatorTok{+}\DecValTok{3}\NormalTok{, cb}\OperatorTok{*}\DecValTok{3}\NormalTok{:cb}\OperatorTok{*}\DecValTok{3}\OperatorTok{+}\DecValTok{3}\NormalTok{] }\OperatorTok{=} \FloatTok{1.0}
        \ControlFlowTok{elif}\NormalTok{ result }\OperatorTok{==} \DecValTok{2}\NormalTok{:}
\NormalTok{            tensor[}\DecValTok{4}\NormalTok{, rb}\OperatorTok{*}\DecValTok{3}\NormalTok{:rb}\OperatorTok{*}\DecValTok{3}\OperatorTok{+}\DecValTok{3}\NormalTok{, cb}\OperatorTok{*}\DecValTok{3}\NormalTok{:cb}\OperatorTok{*}\DecValTok{3}\OperatorTok{+}\DecValTok{3}\NormalTok{] }\OperatorTok{=} \FloatTok{1.0}

        \CommentTok{\# Channel 5: active sub{-}board mask}
\NormalTok{        is\_active }\OperatorTok{=}\NormalTok{ (state.active\_sub\_board }\OperatorTok{==} \OperatorTok{{-}}\DecValTok{1} \KeywordTok{or}\NormalTok{ state.active\_sub\_board }\OperatorTok{==}\NormalTok{ sb)}
        \ControlFlowTok{if}\NormalTok{ is\_active }\KeywordTok{and}\NormalTok{ result }\OperatorTok{==} \DecValTok{0}\NormalTok{:}
\NormalTok{            tensor[}\DecValTok{5}\NormalTok{, rb}\OperatorTok{*}\DecValTok{3}\NormalTok{:rb}\OperatorTok{*}\DecValTok{3}\OperatorTok{+}\DecValTok{3}\NormalTok{, cb}\OperatorTok{*}\DecValTok{3}\NormalTok{:cb}\OperatorTok{*}\DecValTok{3}\OperatorTok{+}\DecValTok{3}\NormalTok{] }\OperatorTok{=} \FloatTok{1.0}

    \CommentTok{\# Channel 6: current player indicator (all 1s for current player, 0s for opponent)}
\NormalTok{    tensor[}\DecValTok{6}\NormalTok{, :, :] }\OperatorTok{=} \FloatTok{1.0} \ControlFlowTok{if}\NormalTok{ p }\OperatorTok{==} \DecValTok{1} \ControlFlowTok{else} \FloatTok{0.0}

    \ControlFlowTok{return}\NormalTok{ tensor}
\end{Highlighting}
\end{Shaded}

The 7 channels encode different types of information:

\begin{itemize}
\tightlist
\item
  \textbf{Ch 0} --- my pieces (current player's marks)
\item
  \textbf{Ch 1} --- opponent's pieces
\item
  \textbf{Ch 2} --- sub-boards I have won (filled with 1s across the
  full 3×3 region)
\item
  \textbf{Ch 3} --- sub-boards opponent has won
\item
  \textbf{Ch 4} --- drawn sub-boards
\item
  \textbf{Ch 5} --- currently active sub-board(s): where I am allowed to
  play
\item
  \textbf{Ch 6} --- whose turn it is (a global binary flag, constant
  across the grid)
\end{itemize}

The player-relative encoding (channels 0 and 1 flip depending on whose
turn it is) means the network learns symmetrically --- it always sees
the board from the perspective of the player to move. Without this, the
network would need to learn separate strategies for X and O.

\hypertarget{putting-it-together-a-random-game}{%
\subsection{7. Putting It Together: A Random
Game}\label{putting-it-together-a-random-game}}

Here is a complete random-agent loop using everything we have built:

\begin{Shaded}
\begin{Highlighting}[]
\ImportTok{import}\NormalTok{ random}

\KeywordTok{def}\NormalTok{ play\_random\_game(seed}\OperatorTok{=}\VariableTok{None}\NormalTok{):}
    \ControlFlowTok{if}\NormalTok{ seed }\KeywordTok{is} \KeywordTok{not} \VariableTok{None}\NormalTok{:}
\NormalTok{        random.seed(seed)}
\NormalTok{    state }\OperatorTok{=}\NormalTok{ GameState()}
\NormalTok{    move\_history }\OperatorTok{=}\NormalTok{ []}

    \ControlFlowTok{while} \KeywordTok{not}\NormalTok{ state.is\_terminal:}
\NormalTok{        legal }\OperatorTok{=}\NormalTok{ get\_legal\_moves(state)}
\NormalTok{        move }\OperatorTok{=}\NormalTok{ random.choice(legal)}
\NormalTok{        move\_history.append(move)}
\NormalTok{        state }\OperatorTok{=}\NormalTok{ apply\_move(state, move)}

    \ControlFlowTok{return}\NormalTok{ state, move\_history}

\NormalTok{state, history }\OperatorTok{=}\NormalTok{ play\_random\_game(seed}\OperatorTok{=}\DecValTok{42}\NormalTok{)}
\BuiltInTok{print}\NormalTok{(}\SpecialStringTok{f"Game ended after }\SpecialCharTok{\{}\NormalTok{state}\SpecialCharTok{.}\NormalTok{move\_count}\SpecialCharTok{\}}\SpecialStringTok{ moves"}\NormalTok{)}
\BuiltInTok{print}\NormalTok{(}\SpecialStringTok{f"Winner: }\SpecialCharTok{\{}\StringTok{\textquotesingle{}X\textquotesingle{}} \ControlFlowTok{if}\NormalTok{ state}\SpecialCharTok{.}\NormalTok{winner }\OperatorTok{==} \DecValTok{1} \ControlFlowTok{else} \StringTok{\textquotesingle{}O\textquotesingle{}} \ControlFlowTok{if}\NormalTok{ state}\SpecialCharTok{.}\NormalTok{winner }\OperatorTok{==} \OperatorTok{{-}}\DecValTok{1} \ControlFlowTok{else} \StringTok{\textquotesingle{}Draw\textquotesingle{}}\SpecialCharTok{\}}\SpecialStringTok{"}\NormalTok{)}
\BuiltInTok{print}\NormalTok{(}\SpecialStringTok{f"Legal moves at each step ranged from }\SpecialCharTok{\{}\BuiltInTok{min}\NormalTok{(}\BuiltInTok{len}\NormalTok{(get\_legal\_moves(GameState())) }\ControlFlowTok{for}\NormalTok{ \_ }\KeywordTok{in}\NormalTok{ [}\DecValTok{0}\NormalTok{])}\SpecialCharTok{\}}\SpecialStringTok{ to 81"}\NormalTok{)}
\end{Highlighting}
\end{Shaded}

With seed 42, the game lasts 42 moves and X wins. Running 1000 random
games gives a sense of the distribution:

\begin{Shaded}
\begin{Highlighting}[]
\NormalTok{results }\OperatorTok{=}\NormalTok{ \{}\StringTok{\textquotesingle{}X\textquotesingle{}}\NormalTok{: }\DecValTok{0}\NormalTok{, }\StringTok{\textquotesingle{}O\textquotesingle{}}\NormalTok{: }\DecValTok{0}\NormalTok{, }\StringTok{\textquotesingle{}Draw\textquotesingle{}}\NormalTok{: }\DecValTok{0}\NormalTok{\}}
\NormalTok{lengths }\OperatorTok{=}\NormalTok{ []}

\ControlFlowTok{for}\NormalTok{ seed }\KeywordTok{in} \BuiltInTok{range}\NormalTok{(}\DecValTok{1000}\NormalTok{):}
\NormalTok{    s, h }\OperatorTok{=}\NormalTok{ play\_random\_game(seed}\OperatorTok{=}\NormalTok{seed)}
\NormalTok{    lengths.append(s.move\_count)}
\NormalTok{    key }\OperatorTok{=} \StringTok{\textquotesingle{}X\textquotesingle{}} \ControlFlowTok{if}\NormalTok{ s.winner }\OperatorTok{==} \DecValTok{1} \ControlFlowTok{else} \StringTok{\textquotesingle{}O\textquotesingle{}} \ControlFlowTok{if}\NormalTok{ s.winner }\OperatorTok{==} \OperatorTok{{-}}\DecValTok{1} \ControlFlowTok{else} \StringTok{\textquotesingle{}Draw\textquotesingle{}}
\NormalTok{    results[key] }\OperatorTok{+=} \DecValTok{1}

\BuiltInTok{print}\NormalTok{(}\SpecialStringTok{f"Results: }\SpecialCharTok{\{}\NormalTok{results}\SpecialCharTok{\}}\SpecialStringTok{"}\NormalTok{)}
\BuiltInTok{print}\NormalTok{(}\SpecialStringTok{f"Avg game length: }\SpecialCharTok{\{}\BuiltInTok{sum}\NormalTok{(lengths)}\OperatorTok{/}\BuiltInTok{len}\NormalTok{(lengths)}\SpecialCharTok{:.1f\}}\SpecialStringTok{ moves"}\NormalTok{)}
\BuiltInTok{print}\NormalTok{(}\SpecialStringTok{f"Min/Max: }\SpecialCharTok{\{}\BuiltInTok{min}\NormalTok{(lengths)}\SpecialCharTok{\}}\SpecialStringTok{/}\SpecialCharTok{\{}\BuiltInTok{max}\NormalTok{(lengths)}\SpecialCharTok{\}}\SpecialStringTok{"}\NormalTok{)}
\end{Highlighting}
\end{Shaded}

Typical output:

\begin{verbatim}
Results: {'X': 498, 'O': 487, 'Draw': 15}
Avg game length: 39.4 moves
Min/Max: 17/81
\end{verbatim}

Random play is nearly symmetric between X and O (first-mover advantage
is small), games typically last around 40 moves, and draws are rare.
This baseline is useful: any trained agent should dramatically
outperform random play, and any evaluation should use both colours.

\hypertarget{what-we-have}{%
\subsection{8. What We Have}\label{what-we-have}}

At this point we have a complete, self-contained UTTT game engine:

\begin{itemize}
\tightlist
\item
  \textbf{\texttt{GameState}} --- immutable board representation
\item
  \textbf{\texttt{encode\_move} / \texttt{decode\_move}} --- flat
  integer move encoding
\item
  \textbf{\texttt{get\_legal\_moves}} --- full rule-compliant move
  generation
\item
  \textbf{\texttt{get\_legal\_move\_mask}} --- binary mask for the
  neural network
\item
  \textbf{\texttt{apply\_move}} --- state transition function
\item
  \textbf{\texttt{check\_sub\_board\_winner} /
  \texttt{check\_meta\_winner}} --- win detection
\item
  \textbf{\texttt{encode\_state}} --- 7-channel tensor encoding
\end{itemize}

This is everything the MCTS and training loop need from the game engine.
P2 will wire up simple RL agents (a tabular Q-learner and a neural
Q-network) against this engine, which will make the limitations of both
immediately apparent and motivate the AlphaZero approach.

\begin{center}\rule{0.5\linewidth}{0.5pt}\end{center}

\emph{Theory: \href{essay1_what_is_uttt.md}{T1 --- What is Ultimate
Tic-Tac-Toe?} covers the rules and strategic structure of the game.
\href{essay1b_rl_background.md}{T2 --- The Reinforcement Learning
Landscape} covers the learning algorithms that P2 will implement.}

\begin{center}\rule{0.5\linewidth}{0.5pt}\end{center}

\emph{Full notebook:
\href{https://colab.research.google.com/github/thltsui/UlltimateTicTacToe/blob/Substack/substack/notebook_1_uttt_engine.ipynb}{Notebook
1 --- The UTTT Game Engine} --- runs all the code in this post with
additional tests, an ASCII board visualiser, and a random-agent
tournament.}
